\section{Is Another Search Worth It?}
\label{sec:results}

Tables~\ref{tab:pandora-study} and~\ref{tab:ucp-pandora-study} report separately
calibrated held-out comparisons across the declared inspection-cost sensitivity grid.
Difference is adaptive policy minus costed search-all, so
negative values favor the adaptive policy. The Shopify interval is product-clustered;
the UCP interval is SKU-clustered, so repeated daily decks never become independent
evidence.

\begin{table}[htbp]
  \centering
  \small
  % Generated by paperkit. Do not edit.
\begin{tabular}{@{}rrrrrl@{}}
\toprule
Scenario cost & Comparator & Difference & 95\% CI & Products & Status \\
\midrule
catalog expansion (\$2.52) & search all & +0.0044 & [-0.0125, 0.0259] & 16 & No reliable improvement \\
agentic review 2m (\$10.21) & search all & -0.0695 & [-0.2234, 0.0897] & 16 & No reliable improvement \\
long horizon research 5m (\$25.55) & search all & -0.5528 & [-0.9130, -0.2373] & 16 & Adaptive improvement \\
\bottomrule
\end{tabular}

  \caption{Held-out costly-search study on a frozen Shopify seller-deck snapshot.
  Each row has four independent held-out product decks. Search-all opens every
  available card and pays every declared inspection cost. Lower total acquisition
  cost is better.}
  \label{tab:pandora-study}
\end{table}

\begin{table}[htbp]
  \centering
  \small
  % Generated by paperkit. Do not edit.
\begin{tabular}{@{}rrrrrl@{}}
\toprule
Scenario cost & Comparator & Difference & 95\% CI & SKU clusters & Status \\
\midrule
catalog expansion (\$2.52) & search all & -0.0847 & [-0.1653, -0.0188] & 1381 & Adaptive improvement \\
agentic review 2m (\$10.21) & search all & -0.7909 & [-1.1296, -0.5236] & 1381 & Adaptive improvement \\
long horizon research 5m (\$25.55) & search all & -2.1914 & [-3.0330, -1.5117] & 1381 & Adaptive improvement \\
\bottomrule
\end{tabular}

  \caption{Held-out costly-search replay on the quality-included September UCP daily
  panel. Each SKU and all of its dates are assigned together to calibration or
  evaluation. Search-all opens every available card and pays every declared inspection
  cost. Merchant domains are observed catalog cards, not verified checkout offers.
  Lower total acquisition cost is better.}
  \label{tab:ucp-pandora-study}
\end{table}

% Claim PANDORA-ADVANTAGE-001.
The UCP daily panel has a held-out interval favoring adaptive stopping in every
registered real-dollar scenario when resampled by SKU. The smaller Shopify replay
favors adaptive stopping in its long-horizon scenario, while its two lower-cost
intervals remain inconclusive. The measured answer is therefore that adaptive stopping
can earn its cost by avoiding unproductive remaining inspections, relative to always
searching every available seller.

Cost sensitivity identifies the expected boundary. Table~\ref{tab:ucp-cost-sensitivity}
holds the dated catalog-expansion provider-rate workload fixed and multiplies its
complete per-inspection cost. At zero through one-tenth of the workload, adaptive and
search-all have three distinct regions. At zero cost, search-all reliably has lower
total cost, as free recall makes opening every feasible card weakly optimal. At
one-hundredth and one-tenth of the workload, the interval crosses zero. At the current
workload and ten times it,
adaptive stopping reliably lowers total cost because it can stop before paying for
low-value remaining inspections.

\begin{table}[htbp]
  \centering
  \small
  % Generated by paperkit. Do not edit.
\begin{tabular}{@{}rrrrrl@{}}
\toprule
Cost multiple & Cost & Comparator & Difference & 95\% CI & Status \\
\midrule
0$\times$ & \$0.00000 & search all & +0.0003 & [0.0001, 0.0007] & Search-all improvement \\
0.01$\times$ & \$0.02523 & search all & -0.0000 & [-0.0006, 0.0007] & No reliable improvement \\
0.1$\times$ & \$0.25228 & search all & -0.0043 & [-0.0119, 0.0012] & No reliable improvement \\
1$\times$ & \$2.52275 & search all & -0.0847 & [-0.1660, -0.0168] & Adaptive improvement \\
10$\times$ & \$25.22750 & search all & -2.1584 & [-2.9865, -1.4523] & Adaptive improvement \\
\bottomrule
\end{tabular}

  \caption{Held-out UCP recall-aware cost sensitivity. Multiples scale the complete
  declared catalog-expansion workload, whose current $1\times$ cost is generated from
  the dated provider rate card. Difference is adaptive minus costed search-all;
  lower is better.}
  \label{tab:ucp-cost-sensitivity}
\end{table}

% Claim UCP-MARKET-RATE-SENSITIVITY-001.

This result exercises calibration-only fitting and clustered evaluation on collected
seller offers and catalog cards. It demonstrates the operational mechanism directly:
continue when a further reveal is worth its price, and stop otherwise. A later study
can broaden the sample and measure deployment-specific costs, but that is extension,
not a prerequisite for reporting this one.